\markdownRendererDocumentBegin
\markdownRendererWarning{The `hybrid` option has been soft-deprecated.}{The `hybrid` option has been soft-deprecated.}{Consider using one of the following better options for mixing TeX and markdown: `contentBlocks`, `rawAttribute`, `texComments`, `texMathDollars`, `texMathSingleBackslash`, and `texMathDoubleBackslash`. For more information, see the user manual at <https://witiko.github.io/markdown/>.}{Consider using one of the following better options for mixing TeX and markdown: `contentBlocks`, `rawAttribute`, `texComments`, `texMathDollars`, `texMathSingleBackslash`, and `texMathDoubleBackslash`. For more information, see the user manual at <https://witiko.github.io/markdown/>.}\markdownRendererSectionBegin
\markdownRendererHeadingOne{Geometria}\markdownRendererInterblockSeparator
{}\markdownRendererSectionBegin
\markdownRendererHeadingTwo{Retornar el angulo mas pequeño entre dos vectores}\markdownRendererInterblockSeparator
{}\markdownRendererInputFencedCode{./_markdown_main/dcb37ac05af8b1c653b75f71823f397c.verbatim}{cpp}{cpp}\markdownRendererInterblockSeparator
{}
\markdownRendererSectionEnd \markdownRendererSectionBegin
\markdownRendererHeadingTwo{Area de la interseccion de dos circulos}\markdownRendererInterblockSeparator
{}\markdownRendererInputFencedCode{./_markdown_main/257506febee6a8189e68057f333a73fa.verbatim}{cpp}{cpp}\markdownRendererInterblockSeparator
{}
\markdownRendererSectionEnd \markdownRendererSectionBegin
\markdownRendererHeadingTwo{Punto de interesccion entre dos circulos}\markdownRendererInterblockSeparator
{}\markdownRendererInputFencedCode{./_markdown_main/9c0b2dee0322dc221b10640963b61540.verbatim}{cpp}{cpp}\markdownRendererInterblockSeparator
{}
\markdownRendererSectionEnd \markdownRendererSectionBegin
\markdownRendererHeadingTwo{Puntos mas cercanos}\markdownRendererInterblockSeparator
{}\markdownRendererInputFencedCode{./_markdown_main/de22892c3ebc696b6c4355724dd26373.verbatim}{cpp}{cpp}\markdownRendererInterblockSeparator
{}
\markdownRendererSectionEnd \markdownRendererSectionBegin
\markdownRendererHeadingTwo{Determinar si 3 puntos son coliniales}\markdownRendererInterblockSeparator
{}Estan en el misma linea\markdownRendererInterblockSeparator
{}\markdownRendererInputFencedCode{./_markdown_main/11e72fad0d4e4e76e112e1aa422656bb.verbatim}{cpp}{cpp}\markdownRendererInterblockSeparator
{}
\markdownRendererSectionEnd \markdownRendererSectionBegin
\markdownRendererHeadingTwo{Determinar si puntos 3D son coplanos}\markdownRendererInterblockSeparator
{}En el mismo plano\markdownRendererInterblockSeparator
{}\markdownRendererInputFencedCode{./_markdown_main/95245d26278febf6548d5f6d06150093.verbatim}{cpp}{cpp}\markdownRendererInterblockSeparator
{}
\markdownRendererSectionEnd \markdownRendererSectionBegin
\markdownRendererHeadingTwo{Intereseccion de una linea con un circulo}\markdownRendererInterblockSeparator
{}\markdownRendererInputFencedCode{./_markdown_main/69e3a166f20e5546919c45e61cdcf252.verbatim}{cpp}{cpp}\markdownRendererInterblockSeparator
{}
\markdownRendererSectionEnd \markdownRendererSectionBegin
\markdownRendererHeadingTwo{Convertir puntos a forma general}\markdownRendererInterblockSeparator
{}ax + b + c = 0\markdownRendererInterblockSeparator
{}\markdownRendererInputFencedCode{./_markdown_main/8f3615b79d53a65d8545effeb6a40a0f.verbatim}{cpp}{cpp}\markdownRendererInterblockSeparator
{}
\markdownRendererSectionEnd \markdownRendererSectionBegin
\markdownRendererHeadingTwo{Distancia entre coordenadas geograficas}\markdownRendererInterblockSeparator
{}\markdownRendererInputFencedCode{./_markdown_main/2b15251824c597c52c0a171b01c013aa.verbatim}{cpp}{cpp}\markdownRendererInterblockSeparator
{}
\markdownRendererSectionEnd \markdownRendererSectionBegin
\markdownRendererHeadingTwo{Area de Triangulo}\markdownRendererInterblockSeparator
{}\markdownRendererInputFencedCode{./_markdown_main/08eed725d6f2febdcbe324ab240604c6.verbatim}{cpp}{cpp}\markdownRendererInterblockSeparator
{}
\markdownRendererSectionEnd \markdownRendererSectionBegin
\markdownRendererHeadingTwo{Interseccion punto con Circulo}\markdownRendererInterblockSeparator
{}Given a circle and a point around or inside the circle we wish to find place(s) of intersection\markdownRendererSoftLineBreak
{}of the lines from the point which are tangent to the circle.\markdownRendererInterblockSeparator
{}\markdownRendererInputFencedCode{./_markdown_main/fca4132f76db8fd3eca890057d45fb26.verbatim}{cpp}{cpp}\markdownRendererInterblockSeparator
{}
\markdownRendererSectionEnd \markdownRendererSectionBegin
\markdownRendererHeadingTwo{Rotar Puntos}\markdownRendererInterblockSeparator
{}\markdownRendererInputFencedCode{./_markdown_main/58b1ea81c1cf6fac315ad41b96b41f2e.verbatim}{cpp}{cpp}
\markdownRendererSectionEnd 
\markdownRendererSectionEnd \markdownRendererDocumentEnd